\pagenumbering{roman} 
\setcounter{page}{1} 
\pagestyle{plain}

%%%%%%%%%%%%%%%%
%%% CREDITOS %%%
%%%%%%%%%%%%%%%%
\chapter*{Créditos/Copyright}

%Una página con la especificación de créditos/copyright para el proyecto (ya sea aplicación por un lado y documentación por el otro, o unificadamente), así como la del uso de marcas, productos o servicios de terceros (incluidos códigos fuente). Si una persona diferente al autor colaboró en el proyecto, tiene que quedar explicitada su identidad y qué hizo.

%A continuación se ejemplifica el caso más habitual, aunque se puede modificar por cualquier otra alternativa:

\vspace{1cm}

\begin{figure}[ht]
    \centering
	\includegraphics[scale=0.2]{images/license.png}
\end{figure}

Esta obra está sujeta a una licencia de Reconocimiento-CompartirIgual \href{https://creativecommons.org/licenses/by-sa/4.0/deed.es}{4.0 Internacional (CC BY-SA 4.0)}.

%%%%%%%%%%%%%
%%% FICHA %%%
%%%%%%%%%%%%%
\chapter*{FICHA DEL TRABAJO FINAL}

\begin{table}[ht]
	\centering{}
	\renewcommand{\arraystretch}{2}
	\begin{tabular}{r | l}
		\hline
		Título del trabajo: &  Modelización de temas de llamadas en tiempo real\\
		\hline
       Nombre del autor: & Manuel E. Gómez Montero\\
		\hline
        Tutor UOC: & Ana Valdivia Garcia\\
		\hline
		Tutor TME: & Antonio Fernández Gallardo\\
		\hline
        Nombre del PRA: & Jordi Casas Roma\\
		\hline
        Fecha de entrega (mm/aaaa): & 01/2019\\
		\hline
        Titulación o programa: & Máster Universitario en Ciencia de Datos\\
		\hline
        Área del Trabajo Final: & Procesamiento del Lenguaje Natural\\
		\hline
        Idioma del trabajo: & Español\\
		\hline
        Palabras clave & ``natural language processing'', ``sentiment analysis'',\\
        & ``real time'', ``call center'', ``topic modeling'', ``deep learning''\\
		\hline
	\end{tabular}
\end{table}

%%%%%%%%%%%%%%%%%%%
%%% DEDICATORIA %%%
%%%%%%%%%%%%%%%%%%%
%\chapter*{Dedicatoria/Cita}

%Breves palabras de dedicatoria y/o una cita.
\cleardoublepage
\begin{dedication}
A Ana, por el tiempo robado
\end{dedication}
\cleardoublepage
%%%%%%%%%%%%%%%%%%%
%%% Agradecimientos %%%
%%%%%%%%%%%%%%%%%%%
\chapter*{Agradecimientos}


Un \textit{handicap} a la hora de realizar el proyecto dentro de una gran empresa ha sido el hecho de trabajar con unos plazos tan ajustados. Aspectos como la autorización en el acceso a la información, el acceso a diferentes entornos, la intercomunicación entre áreas, etc. requieren unos tiempos que pueden retrasar la ejecución de un trabajo de este tipo. 
 
Este apartado tiene como objetivo agradecer a todas las personas que han hecho posible la realización de este proyecto en tiempo y forma.


A Willy Gavilán y Carolina Bouvard por hacer posible la  realización de este proyecto dentro de Telefónica. A Antonio Fernández por aceptar que realice este trabajo con su equipo y por tutorizarlo. A Jorge Ayuso y Rus Mesas por su ayuda a la hora de obtener los datos, por sus consejos y por permitirme usar con ellos el equipo con GPUs usados para  el entrenamiento y optimización de modelos. A Laura Canga por permitirme usar el Bus Apache Kafka y a José Ramón Fernández por ayudarme, incluso a deshoras, con la creación de \textit{topics} y usuarios del bus.  A Pablo Palomares por su ayuda a la hora de acceder y adaptar a nuestras necesidades las plataformas \textit{DevOps}: Bitbucket, Nexus y Jenkins. A Nacho Charfolé por darme su visión más allá de la parte técnica y por hacerme entender las ventajas que pueden aportar este tipo de proyectos a una empresa. Y por último, pero no menos importante, a mi tutora, Ana Valdivia, por sus correcciones minuciosas y por guiarme y animarme en todo momento.

Sin todos vosotros un trabajo de este tipo no sería posible. \textbf{{\LARGE Gracias}}.


%%%%%%%%%%%%%%%%
%%% RESUMEN  %%%
%%%%%%%%%%%%%%%%
\chapter*{Resumen}
\addcontentsline{toc}{chapter}{Resumen}


Un \textit{call center} es el área de una empresa el cuál se encarga de recibir y transmitir llamadas desde o hacia clientes, socios comerciales u otras compañías externas. La gran cantidad de información generada a partir de estas llamada refleja, entre otros aspectos, las necesidades de los clientes de la empresa. 


La capacidad de clasificar las llamadas recibidas en tiempo real en función de la temática , provocaría en las empresas una mejora en la calidad de atención al público y en la experiencia del cliente.  Esto es abordable mediante técnicas de \textit{machine learning}.

%Por ello, resulta una tarea esencial ser capaces de reaccionar en el menor tiempo posible ante un cambio en la temática de las llamadas con el objetivo de mejorar la experiencia del cliente.

%optimizar el tiempo de respuesta ante un cambio en la temática de las llamadas, para así reaccionar en tiempo real a las peticiones de los mismos y mejorar la percepción que estos tienen sobre la compañía. 

%Las empresas mejorarían la calidad de las llamadas de atención al público si fuesen capaces de clasificar en tiempo real las llamadas que reciben o emiten, en función de la temática que se este tratando en cada momento. Esto es abordable mediante técnicas de machine learning. 

En este documento se presenta un sistema implementado en el \textit{call center} de Telefónica que nos 
 permite, a partir de la transcripción de las llamadas, obtener la temática de las mismas utilizando métodos de Procesamiento de Lenguaje Natural. El análisis expuesto realiza clasificaciones con diferentes arquitecturas, tanto con modelos no supervisados como con modelos supervisados de aprendizaje profundo. Dependiendo del modelo se han estudiado también diferentes métodos a la hora de representar las transcripciones.
 
 Por último, se ha llevado uno de los modelos a producción, dotando al sistema de la posibilidad de realizar esta clasificación en tiempo real y proporcionando a la empresa la capacidad de visualizar la evolución de las llamadas a lo largo del tiempo y de generar alertas en base a anomalías. El sistema llevado a producción ha sido totalmente desplegado en contenedores y provisto de mecanismos  de monitorización, integración continua y despliegue continuo. 
 
El proyecto presentado abre a la empresa un gran número de posibilidades, tanto a la hora de  realizar nuevos sistemas de Procesamiento del Lenguaje Natural, como a la hora de llevar modelos de aprendizaje profundo a un entorno productivo.
 





%El documento que se presenta tiene como objetivo final mejorar la operatividad del \textit{call-center}  de una gran empresa, extrayendo mediante técnicas de \textit{machine learning} la temática de las llamadas que se realizan al mismo y dando la posibilidad a la empresa de reaccionar en tiempo real, en función de la temática que se este tratando en cada momento. 

%En este documento se presenta un sistema implementado en el \textit{call-center} de Telefónica que nos  permite, a partir de la transcripción de las llamadas, descubrir en tiempo real la temática de las mismas. Esta modelización de \textit{topics} se ha realizado utilizando métodos de Procesamiento de Lenguaje Natural y aprendizaje profundo. El sistema realiza la clasificación de las nuevas llamadas en tiempo real, permitiendo a los usuarios visualizar la evolución en la temática de las mismas y generar alertas en base a anomalías.  


\chapter*{Abstract}
\addcontentsline{toc}{chapter}{Abstract}

A call center is the area of a company in charge of receiving and transmitting calls to or from customers, business partners or other external companies. The vast amount of information generated from these calls reflects, among other aspects, the needs of the company's customers.


The ability to classify the calls received in real time according to the subject, would lead to an improvement in the quality of the company´s customer service as well as the customer experience. This is approachable through machine learning techniques.


This draft reveals a system implemented in the call center of Telefónica that allows us, through the transcription of the calls, to know the topic of the calls using Natural Language Processing methods. The analysis mentioned above makes classifications with different architectures, with unsupervised models as well as deep learning supervised models. Depending on the model, different methods have also been studied when representing the transcripts.


Finally, one of these models has been moved to production, giving the system the possibility of making this classification in real time and providing the company the ability of visualizing the evolution of the calls over time and generating alerts in the event of any anomaly. The system moved to production has been fully deployed in containers and equipped with monitoring mechanisms, continuous integration and continuous deployment.


The project opens up a vast number of possibilities for the company to carry out new Natural Language Processing systems. Also, to bring deep learning models to a productive environment. 




\onehalfspacing


